%% Signal porte Pi_theta (largeur theta, hauteur 1, centre en 0) et sa
%% TF, theta*sinc(pi f theta) (theta=1 ici). Exemple fil rouge de tout
%% le chapitre. Repris du poly (Exemple, section 2). Consomme par
%% slides.
\begin{tikzpicture}[x=1.2cm,y=1cm]
    \draw[-Latex] (-2.2,0) -- (2.2,0) node[right] {$t$};
    \draw[-Latex] (0,-0.2) -- (0,1.5) node[above] {$\Pi_\theta(t)$};
    \draw[niceblue,thick] (-2,0) -- (-1,0) -- (-1,1) -- (1,1) -- (1,0) -- (2,0);
    \draw[dashed] (-1,0) -- (-1,1); \draw[dashed] (1,0) -- (1,1);
    \node[below] at (-1,0) {$-\theta/2$};
    \node[below] at (1,0) {$\theta/2$};
    \node[left] at (0,1) {$1$};
\end{tikzpicture}
\hspace{0.3cm}
{\LARGE $\overunderset{\mathcal F}{\mathcal{F}^{-1}}{\rightleftarrows}$}
\hspace{0.3cm}
\begin{tikzpicture}
    \begin{axis}[
        width=6.5cm, height=4cm,
        xlabel={$f$}, ylabel={$X(f)$},
        xmin=-6, xmax=6,
        axis lines=middle,
        xlabel style={at=(current axis.right of origin), anchor=west},
        ylabel style={at=(current axis.above origin), anchor=south},
    ]
    \addplot[domain=-6:6, samples=200, niceblue, thick] {sin(deg(3.14159*x))/(3.14159*x)};
    \end{axis}
\end{tikzpicture}
